This framework is written for the clinicians who will be asked to supervise
autonomous Physical AI in oncology clinical trials, and who must decide, at the
bedside, when that reliance is earned. It rests on a single mechanism, verification
before generation, in which a software agent proposes a clinical action and a
ten-gate verification, validation, and uncertainty quantification examination
either accepts it under audit, escalates it to a qualified human, or blocks it
before it can reach a patient. The decision to trust is organized as eight plain
questions a clinician already asks of any tool: is it competent, is it safe, is it
transparent, am I in control, is it equitable, is it reliable, who is accountable,
and does it fit my workflow. Each question pairs a clinical concern with a
credible, cited fact, and closes on the concrete mechanism, a gate, a record, a
role, or a control, that answers it. The aim is not maximal trust but calibrated
trust: reliance proportioned to demonstrated capability, avoiding both automation
bias and algorithm aversion. The framework draws on a documented lineage, including a surgical-humanoid simulation that passed 172 of 172
automated tests across a ten-gate suite, and on the literature of trust in
automation, clinical decision support, and algorithmic bias. Three questions recur
throughout: would I let this system near my own family member, can I defend this
decision at tumor board, and could I show an auditor exactly what happened and why.
The conclusion is an adoption decision a clinician can sign.
